\section{Introdução}
\label{sec:introducao}

O Problema da Ponte e da Tocha é um quebra-cabeça clássico utilizado no
estudo de algoritmos de busca em Inteligência Artificial. Um grupo de
pessoas precisa atravessar uma ponte estreita durante a noite, dispondo de
apenas uma tocha para iluminar o caminho. A ponte suporta, no máximo, duas
pessoas por vez, e qualquer travessia -- de ida ou de volta -- exige que a
tocha esteja com quem está cruzando. Cada pessoa possui um tempo diferente
de travessia e, quando duas pessoas atravessam juntas, o tempo gasto é
determinado pela mais lenta delas. O objetivo é encontrar uma sequência de
travessias que leve todas as pessoas ao lado de destino da ponte no menor
tempo total possível.

Este trabalho tem como objetivo modelar o problema como um Problema em
Espaço de Estados e resolvê-lo por meio de diferentes algoritmos de busca,
comparando suas garantias teóricas e seu comportamento prático. Foram
implementados dois algoritmos de busca não informada -- Busca em
Profundidade (DFS) e Busca em Largura (BFS) -- e dois algoritmos que
utilizam o custo das transições para orientar a exploração do espaço de
estados -- Busca de Custo Mínimo (UCS) e Busca A*.

O restante do relatório está organizado da seguinte forma: a
Seção~\ref{sec:modelagem} apresenta a modelagem formal do problema como um
espaço de estados; as Seções~\ref{sec:busca-nao-informada}
e~\ref{sec:busca-informada} descrevem a implementação dos quatro algoritmos
de busca; a Seção~\ref{sec:comparacao-algoritmos} apresenta a comparação
experimental entre eles; e a última seção conclui o trabalho, retomando os
principais resultados obtidos.
